\documentclass[10pt,a4paper]{amsart}
\usepackage{amsmath,amssymb,amsthm,graphicx}
\usepackage{verbatim}
\newenvironment{reference}{\comment}{\endcomment}
%\newenvironment{reference}{}{}
\newenvironment{slogan}{\comment}{\endcomment}
\newenvironment{history}{\comment}{\endcomment}

% For commutative diagrams we use Xy-pic
\usepackage[all]{xy}

% We use 2cell for 2-commutative diagrams.
\xyoption{2cell}
% Theorem environments.
%
\theoremstyle{plain}
\newtheorem{theorem}[subsection]{Theorem}
\newtheorem{proposition}[subsection]{Proposition}
\newtheorem{lemma}[subsection]{Lemma}

\theoremstyle{definition}
\newtheorem{definition}[subsection]{Definition}
\newtheorem{example}[subsection]{Example}
\newtheorem{exercise}[subsection]{Exercise}
\newtheorem{situation}[subsection]{Situation}

\theoremstyle{remark}
\newtheorem{remark}[subsection]{Remark}
\newtheorem{remarks}[subsection]{Remarks}

\numberwithin{equation}{subsection}

% Macros
%
\def\lim{\mathop{\mathrm{lim}}\nolimits}
\def\colim{\mathop{\mathrm{colim}}\nolimits}
\def\Spec{\mathop{\mathrm{Spec}}}
\def\Hom{\mathop{\mathrm{Hom}}\nolimits}
\def\Ext{\mathop{\mathrm{Ext}}\nolimits}
\def\SheafHom{\mathop{\mathcal{H}\!\mathit{om}}\nolimits}
\def\SheafExt{\mathop{\mathcal{E}\!\mathit{xt}}\nolimits}
\def\Sch{\mathit{Sch}}
\def\Mor{\mathop{\mathrm{Mor}}\nolimits}
\def\Ob{\mathop{\mathrm{Ob}}\nolimits}
\def\Sh{\mathop{\mathit{Sh}}\nolimits}
\def\NL{\mathop{N\!L}\nolimits}
\def\CH{\mathop{\mathrm{CH}}\nolimits}
\def\proetale{{pro\text{-}\acute{e}tale}}
\def\etale{{\acute{e}tale}}
\def\QCoh{\mathit{QCoh}}
\def\Ker{\mathop{\mathrm{Ker}}}
\def\Im{\mathop{\mathrm{Im}}}
\def\Coker{\mathop{\mathrm{Coker}}}
\def\Coim{\mathop{\mathrm{Coim}}}

% Boxtimes
%
\DeclareMathSymbol{\boxtimes}{\mathbin}{AMSa}{"02}

%
% Macros for moduli stacks/spaces
%
\def\QCohstack{\mathcal{QC}\!\mathit{oh}}
\def\Cohstack{\mathcal{C}\!\mathit{oh}}
\def\Spacesstack{\mathcal{S}\!\mathit{paces}}
\def\Quotfunctor{\mathrm{Quot}}
\def\Hilbfunctor{\mathrm{Hilb}}
\def\Curvesstack{\mathcal{C}\!\mathit{urves}}
\def\Polarizedstack{\mathcal{P}\!\mathit{olarized}}
\def\Complexesstack{\mathcal{C}\!\mathit{omplexes}}
% \Pic is the operator that assigns to X its picard group, usage \Pic(X)
% \Picardstack_{X/B} denotes the Picard stack of X over B
% \Picardfunctor_{X/B} denotes the Picard functor of X over B
\def\Pic{\mathop{\mathrm{Pic}}\nolimits}
\def\Picardstack{\mathcal{P}\!\mathit{ic}}
\def\Picardfunctor{\mathrm{Pic}}
\def\Deformationcategory{\mathcal{D}\!\mathit{ef}}
\usepackage{cleveref}
\crefname{originaltheorem}{original theorem}{original theorems}
\Crefname{originaltheorem}{Original theorem}{Original theorems}
\makeatletter
\newlabel{theorem-gabriel-popescu}{{0F5U}{source}{}{sourcepointer-1}{}}
\newlabel{theorem-gabriel-popescu@cref}{{[originaltheorem][0][]0F5U}{[1][0][]source}{}{}{}}
\makeatother
\crefname{originaltheorem}{original theorem}{original theorems}
\Crefname{originaltheorem}{Original theorem}{Original theorems}
\makeatletter
\newlabel{theorem-K-injective-embedding-grothendieck}{{079P}{source}{}{sourcepointer-2}{}}
\newlabel{theorem-K-injective-embedding-grothendieck@cref}{{[originaltheorem][0][]079P}{[1][0][]source}{}{}{}}
\makeatother
\crefname{originallemma}{original lemma}{original lemmas}
\Crefname{originallemma}{Original lemma}{Original lemmas}
\makeatletter
\newlabel{derived-lemma-K-injective-defined}{{070Y}{source}{}{sourcepointer-3}{}}
\newlabel{derived-lemma-K-injective-defined@cref}{{[originallemma][0][]070Y}{[1][0][]source}{}{}{}}
\makeatother
\crefname{originallemma}{original lemma}{original lemmas}
\Crefname{originallemma}{Original lemma}{Original lemmas}
\makeatletter
\newlabel{derived-lemma-right-derived-exact-functor}{{015F}{source}{}{sourcepointer-4}{}}
\newlabel{derived-lemma-right-derived-exact-functor@cref}{{[originallemma][0][]015F}{[1][0][]source}{}{}{}}
\makeatother
\crefname{originallemma}{original lemma}{original lemmas}
\Crefname{originallemma}{Original lemma}{Original lemmas}
\makeatletter
\newlabel{derived-lemma-derived-adjoint-functors}{{09T5}{source}{}{sourcepointer-5}{}}
\newlabel{derived-lemma-derived-adjoint-functors@cref}{{[originallemma][0][]09T5}{[1][0][]source}{}{}{}}
\makeatother
\crefname{originallemma}{original lemma}{original lemmas}
\Crefname{originallemma}{Original lemma}{Original lemmas}
\makeatletter
\newlabel{categories-lemma-adjoint-fully-faithful}{{07RB}{source}{}{sourcepointer-6}{}}
\newlabel{categories-lemma-adjoint-fully-faithful@cref}{{[originallemma][0][]07RB}{[1][0][]source}{}{}{}}
\makeatother
\title{The Gabriel-Popescu adjunction is an equivalence on the module side}
\author{The Stacks Project Authors}
\date{stacks-project:injectives.tex:local-snapshot-20261012; GNU FDL 1.2 or later}
\begin{document}
\maketitle
\noindent\textit{Editorial note.} All delivered mathematical text is the author's own, copied from the pinned Stacks Project chapter file sites-cohomology.tex; no statement or proof step was rewritten, reordered or removed.\par\medskip
\noindent\textit{Editorial note.} The proof is the author's complete argument; the statement carries the author's reference environment recording the provenance of the result, which is a comment-type environment of the chapter preamble so its citation is discarded on delivery, and the argument itself is written out in the counted span.\par\medskip
\noindent\textit{Editorial note.} Every labelled reference occurring inside the counted spans is carried in assembly.external\_references as a non-load-bearing navigation pointer (same-chapter ones at their definition site, cross-chapter ones at the invocation inside the counted span), so no load-bearing mathematics is replaced by a stub. No bibliography citation occurs inside the counted spans.\par\medskip
\noindent\textit{Editorial note.} The same-name check for this card scanned every staged Stacks chapter of this repository -- all sources/stacks-*/pinned-original-source.tex files, the card's own chapter excluded -- and the labels found are recorded in dedup.same\_name\_cross\_chapter on the card.\par\medskip
\noindent\textit{Editorial note.} The preamble is the chapter's own retained minimal preamble (sources/stacks-sites-cohomology/retained-preamble.tex), a verbatim excerpt of the native preamble: the Xy-pic setup and the native theorem declarations with the mathematical macros.\par\medskip
\noindent\textit{Allowed standard prerequisites (editorial).}\begin{itemize}
\item The adjunction of the two functors. Same chapter, the construction of the functors F and G and their adjunction established earlier. Supplies the pair of functors whose derived functors are compared and the identity of the composite on the module side.
\item Fully faithful adjunctions. Categories, Lemma categories-lemma-adjoint-fully-faithful. Supplies that an adjunction whose unit or counit is an isomorphism is fully faithful, which is how full faithfulness of the derived functor is concluded.
\item K-injective representatives. Same chapter, Theorem theorem-K-injective-embedding-grothendieck. Supplies that the derived functor can be computed on a K-injective representative, which is what makes the composite computation available.
\end{itemize}\medskip
\ifdefined\hypertarget\hypertarget{sourcepointer-1}{}\fi
\noindent\textit{Source reference.} Original theorem 0F5U: sources/stacks-injectives/pinned-original-source.tex, source lines 2878--2878. Navigation pointer invoked at proof line 2928; the counted span contains the pointer only and the cited result is not reproduced inside it, so the pointer is registered as a non-load-bearing navigation stub in the established form for Stacks material. No bibliography citation is involved.\par\medskip
\ifdefined\hypertarget\hypertarget{sourcepointer-2}{}\fi
\noindent\textit{Source reference.} Original theorem 079P: sources/stacks-injectives/pinned-original-source.tex, source lines 2024--2024. Navigation pointer invoked at proof line 2942; the counted span contains the pointer only and the cited result is not reproduced inside it, so the pointer is registered as a non-load-bearing navigation stub in the established form for Stacks material. No bibliography citation is involved.\par\medskip
\ifdefined\hypertarget\hypertarget{sourcepointer-3}{}\fi
\noindent\textit{Source reference.} Original lemma 070Y: sources/stacks-derived/pinned-original-source.tex, source lines 9964--9964. Navigation pointer invoked at proof line 2944; the counted span contains the pointer only and the cited result is not reproduced inside it, so the pointer is registered as a non-load-bearing navigation stub in the established form for Stacks material. No bibliography citation is involved.\par\medskip
\ifdefined\hypertarget\hypertarget{sourcepointer-4}{}\fi
\noindent\textit{Source reference.} Original lemma 015F: sources/stacks-derived/pinned-original-source.tex, source lines 6101--6101. Navigation pointer invoked at proof line 2945; the counted span contains the pointer only and the cited result is not reproduced inside it, so the pointer is registered as a non-load-bearing navigation stub in the established form for Stacks material. No bibliography citation is involved.\par\medskip
\ifdefined\hypertarget\hypertarget{sourcepointer-5}{}\fi
\noindent\textit{Source reference.} Original lemma 09T5: sources/stacks-derived/pinned-original-source.tex, source lines 9813--9813. Navigation pointer invoked at proof line 2946; the counted span contains the pointer only and the cited result is not reproduced inside it, so the pointer is registered as a non-load-bearing navigation stub in the established form for Stacks material. No bibliography citation is involved.\par\medskip
\ifdefined\hypertarget\hypertarget{sourcepointer-6}{}\fi
\noindent\textit{Source reference.} Original lemma 07RB: sources/stacks-categories/pinned-original-source.tex, source lines 3657--3657. Navigation pointer invoked at proof line 2893; the counted span contains the pointer only and the cited result is not reproduced inside it, so the pointer is registered as a non-load-bearing navigation stub in the established form for Stacks material. No bibliography citation is involved.\par\medskip
\begin{lemma}
\label{lemma-gabriel-popescu}
\begin{reference}
\cite[Corollary 4.1]{serpe}
\end{reference}
Let $\mathcal{A}$ be a Grothendieck abelian category. Let
$R$, $F$, $G$ be as in the Gabriel-Popescu theorem
(Theorem \ref{theorem-gabriel-popescu}). Then we obtain
derived functors
$$
RG : D(\mathcal{A}) \to D(\text{Mod}_R)
\quad\text{and}\quad
F : D(\text{Mod}_R) \to D(\mathcal{A})
$$
such that $F$ is left adjoint to $RG$, $RG$ is fully faithful,
and $F \circ RG = \text{id}$.
\end{lemma}

\begin{proof}
The existence and adjointness of the functors follows from
Theorems \ref{theorem-gabriel-popescu} and
\ref{theorem-K-injective-embedding-grothendieck}
and
Derived Categories, Lemmas \ref{derived-lemma-K-injective-defined},
\ref{derived-lemma-right-derived-exact-functor}, and
\ref{derived-lemma-derived-adjoint-functors}.
The statement $F \circ RG = \text{id}$ follows because we can
compute $RG$ on an object of $D(\mathcal{A})$ by applying $G$
to a suitable representative complex $I^\bullet$ (for example
a K-injective one) and then $F(G(I^\bullet)) = I^\bullet$
because $F \circ G = \text{id}$. Fully faithfulness of $RG$
follows from this by
Categories, Lemma \ref{categories-lemma-adjoint-fully-faithful}.
\end{proof}

\end{document}
